\SourceLeaf{233}%
reserve for another occasion a fuller account of this subject for figures taken in the greatest generality.

\Article{7}
Let us now seek a formula for the measure of curvature at any point of a curved surface. Let \(\mathrm{d}\sigma\) denote the area of an element of this surface. Then \(Z\mathrm{d}\sigma\) will be the area of its projection onto the plane of the coordinates \(x,y\). Likewise, if \(\mathrm{d}\Sigma\) is the area of the corresponding element on the sphere, \(Z\mathrm{d}\Sigma\) will be the area of its projection onto the same plane. The positive or negative sign of \(Z\) indicates whether the position of the projection is similar or opposite to that of the projected element. Clearly, then, these projections bear the same ratio in magnitude and the same relation in position to one another as the elements themselves. Let us now consider a triangular element on the curved surface and suppose that the coordinates of the three points forming its projection are
\[
\begin{array}{ll}
x,&y\\
x+\mathrm{d}x,&y+\mathrm{d}y\\
x+\delta x,&y+\delta y
\end{array}
\]
Twice the area of this triangle is expressed by the formula
\[
\mathrm{d}x\mathord{.}\delta y-\mathrm{d}y\mathord{.}\delta x
\]
with positive or negative sign according to whether the side from the first point to the third lies relative to the side from the first point to the second as the coordinate axis \(y\) lies relative to the coordinate axis \(x\), or in the opposite way.

Likewise, if the coordinates of the three points forming the projection of the corresponding element on the sphere, with the centre of the sphere as origin, are
\[
\begin{array}{ll}
X,&Y\\
X+\mathrm{d}X,&Y+\mathrm{d}Y\\
X+\delta X,&Y+\delta Y
\end{array}
\]
twice the area of this projection is expressed by
\[
\mathrm{d}X\mathord{.}\delta Y-\mathrm{d}Y\mathord{.}\delta X
\]
and the same rules as above apply to the sign of this expression. Therefore, at this point of the curved surface, %
